\section{Hur ska jag förbättra mina kunskaper och färdigheter}
\label{sec:förbättra mind kunskaper och färdigheter} 

Under LIA-perioden har jag kommit en bra bit på vägen mot min målsättning att bli DevOps Engineer, men jag känner också att jag har mycket kvar att utveckla. Mina arbetsuppgifter handlade mest om Kubernetes och GKE, och därför är detta det område som jag vill prioritera mest framöver. För att kunna arbeta mer självständigt med GKE behöver jag bli ännu säkrare på hur kluster fungerar, hur resurser skapas och hur olika komponenter pratar med varandra.

Samtidigt vet jag att för att arbeta effektivt med Kubernetes i molnet måste jag också ha bra kunskaper i Terraform, eftersom all infrastruktur i Insighta byggs med infrastruktur som kod. Därför vill jag fortsätta öva på Terraform och CDKTF och skapa fler resurser i testmiljöer för att förstå verktyget ännu bättre.

När jag har stärkt mina kunskaper inom Terraform och GKE vill jag gå vidare till nästa steg, som är CI/CD och branch deployment. Detta är något jag redan har börjat arbeta med i slutet av LIA:n, men jag vill fördjupa mig mycket mer så att jag kan bygga stabila och automatiserade pipelines. Det är en viktig del av DevOps-rollen och något jag vill bli riktigt bra på.

Jag vill också förbättra min kunskap inom Dagster, eftersom det är en stor del av Insightas arbetsflöde. Även om Dagster inte är ett klassiskt DevOps-verktyg är det viktigt att jag förstår hur det fungerar när jag ska drifta och övervaka det i Kubernetes. Detta kommer också hjälpa mig senare om jag går vidare mot MLOps.

För att utvecklas planerar jag att:

Fortsätta öva på Kubernetes och GKE under min andra LIA.

Göra egna små projekt där jag bygger kluster och testar olika deployment-sätt.

Fördjupa mig i Terraform/CDKTF och läsa mer dokumentation.

Lära mig mer om CI/CD genom att skapa egna pipelines i GitHub Actions.

Fortsätta arbeta med Dagster och förstå hur det kopplas till GCP och Kubernetes.

Jag tänker börja med detta omedelbart under LIA 2 och under resten av utbildningen. Jag känner att jag är på rätt väg mot att bli DevOps Engineer, men jag behöver fortsätta utvecklas för att nå hela vägen.